% ===========================================================
% Chapter 17: Open Problems and Emerging Directions
% Part IV: Frontiers
% ===========================================================

\chapter{Открытые задачи и перспективные направления}
\label{ch:open}

\begin{quote}
\itshape
Лучший способ предсказать будущее~--- изобрести его.
\par\medskip
\upshape\raggedleft
---\,Alan Kay
\end{quote}

\bigskip

В предыдущих главах обучение с подкреплением было прослежено от
многоруких бандитов (Глава~\ref{ch:bandits}) через динамическое
программирование (Глава~\ref{ch:dp}), методы временных разностей
(Глава~\ref{ch:td}), приближение глубокой функции ценности
(Глава~\ref{ch:value-approx}), градиенты стратегии и алгоритмы актор-критик
(Главы~\ref{ch:pg}--\ref{ch:actor-critic}) вплоть до
современного пересечения RL с большими языковыми моделями
(Главы~\ref{ch:lm-post-training}--\ref{ch:agents}). По пути
мы видели, как RL решает игры, выравнивает языковые модели с предпочтениями людей,
обеспечивает сложные рассуждения и координирует мультиагентные системы.

Тем не менее на каждую решённую RL задачу приходится дюжина открытых. Некоторые
давние\,---\,эффективность по образцам была предметом беспокойства с
первых табличных алгоритмов,\,---\,а другие возникли лишь
недавно, движимые быстрым принятием RL как метода дообучения
для базовых моделей. Эта финальная глава обозревает наиболее
важные открытые вопросы, группирует их в связные темы и
пытается наметить путь вперёд. Она неизбежно более умозрительна,
чем предшествующие главы: цель~--- не доказывать теоремы,
а ориентировать читателя на задачи, которые будут определять следующее
десятилетие исследований.


% ===========================================================
\section{Эффективность по образцам}
\label{sec:open-sample}
% ===========================================================

\index{sample efficiency}
Обучение с подкреплением печально известно своей жаждой данных. Безмодельные
алгоритмы, такие как PPO (Глава~\ref{ch:actor-critic})
и DQN (Глава~\ref{ch:value-approx}), часто требуют миллионов или
даже миллиардов взаимодействий со средой для обучения поведению, которое
ребёнок осваивает за минуты. Разрыв между человеческой и машинной
эффективностью по образцам~--- одна из старейших открытых задач области, и
она остаётся одной из наиболее значимых.

\subsection{Модельно-ориентированный RL и модели мира}
\index{model-based reinforcement learning}
\index{world model}

Наиболее естественный путь к эффективности по образцам~--- \emph{обучить модель
среды} и планировать внутри неё. Если агент может предсказать, что
произойдёт до того, как он действует, каждое реальное взаимодействие становится значительно более
информативным. Ранние модельно-ориентированные методы учили динамику переходов в
табличной или линейной форме (Глава~\ref{ch:dp} обсуждает
аспект планирования); современные подходы обучают \emph{модели мира}\,---\,нейронные
сети, предсказывающие будущие наблюдения, вознаграждения и сигналы завершения.

Статья Ha и Schmidhuber «World Models» (2018) показала, что агент
может обучиться сжатому скрытому представлению своих
наблюдений, а затем обучить стратегию целиком внутри обученной
модели. Hafner et al.\ расширили это направление семейством алгоритмов Dreamer
(Dreamer, DreamerV2, DreamerV3), которые обучают
рекуррентную модель пространства состояний и оптимизируют стратегию путём
обратного распространения через воображаемые траектории в скрытом пространстве.
DreamerV3 достиг производительности уровня человека на десятках непрерывных и
дискретных задач управления\,---\,и даже добывал алмазы в Minecraft,\,---\,используя
примерно в $10$--$100\times$ меньше шагов среды, чем сопоставимые
безмодельные методы.

\index{Dreamer}
\index{MuZero}
MuZero (Schrittwieser et al., 2020) применил другой подход:
вместо предсказания сырых наблюдений он обучал скрытую модель динамики,
чьи скрытые состояния обучались сквозным образом для поддержки планирования
через поиск по дереву Монте-Карло. MuZero достиг производительности
AlphaZero в Го, шахматах и сёги, также освоив игры Atari,
демонстрируя, что планирование в скрытом пространстве может соответствовать или превосходить
безмодельные методы даже в визуально сложных областях.

Несмотря на эти успехи, модельно-ориентированный RL сталкивается с постоянными вызовами.
Накопление ошибок модели\,---\,небольшие неточности в обученной
динамике, накапливающиеся на длинных горизонтах планирования,\,---\,могут производить
стратегии, оптимальные в модели, но катастрофические в реальной
среде. Решение \emph{когда} доверять модели и когда
прибегать к реальным данным (так называемый компромисс Дина; Sutton, 1991)
остаётся в значительной мере эвристическим.

\subsection{Перенос знаний и few-shot RL}
\index{transfer learning!in RL}
\index{few-shot RL}

Дополнительный путь к эффективности по образцам~--- \emph{перенос}:
использование знаний из предыдущих задач для ускорения обучения на новых.
Алгоритмы мета-RL, такие как MAML (Finn et al., 2017) и
RL$^2$ (Duan et al., 2016), обучают инициализацию или алгоритм
внутреннего цикла обучения, который может адаптироваться к новой задаче за горсть
эпизодов. Более недавние подходы на основе базовых моделей рассматривают траектории RL
как последовательности и используют обучение в контексте\,---\,тот же
механизм, что позволяет few-shot промптинг в языковых моделях,\,---\,для
адаптации к новым средам во время вывода без обновлений градиента
(см. Раздел~\ref{sec:open-offline} о Decision Transformer и
связанных методах).

Видение единого универсального агента RL, переносящего знания
через разнообразные области\,---\,«базовой модели RL»,\,---\,остаётся
в значительной мере нереализованным. Gato (Reed et al., 2022) показал, что единственный
трансформер может играть в Atari, создавать подписи к изображениям и управлять роботизированной
рукой, но его производительность по отдельным задачам отставала от специализированных агентов.
Устранение этого разрыва~--- активная область исследований.


% ===========================================================
\section{Спецификация вознаграждения и выравнивание}
\label{sec:open-alignment}
% ===========================================================

\index{reward hacking}
\index{Goodhart's law}
\index{alignment}
Каждый алгоритм RL оптимизирует сигнал вознаграждения. Когда этот сигнал
точно отражает то, что мы хотим, оптимизация полезна; когда нет,
агент использует разрыв\,---\,явление, известное как
\emph{взлом вознаграждения}. Задача~--- экземпляр закона Гудхарта\index{Goodhart's law}: «Когда мера становится целью, она
перестаёт быть хорошей мерой».

Глава~\ref{ch:reward-models} ввела модели вознаграждения, обученные на
сравнениях людей, а Глава~\ref{ch:dpo} показала, как DPO обходит
явное моделирование вознаграждения. Оба подхода смягчают взлом вознаграждения в
какой-то мере, но ни один не устраняет его. По мере того как модели становятся более
мощными, ставки неправильно заданных вознаграждений растут.

\subsection{Масштабируемый надзор}
\index{scalable oversight}

Когда возможности агента превышают возможности оценщика, обратная связь от людей
становится ненадёжной. Модель кодирования может производить решения, верные, но непостижимые;
модель научного рассуждения может предлагать эксперименты, чью обоснованность
не может судить неспециалист. Задача \emph{масштабируемого надзора} спрашивает:
как обеспечить надёжный обучающий сигнал для систем, чьи выходы нам нелегко проверить?

Несколько исследовательских программ решают этот вопрос. \emph{Дебаты}
(Irving et al., 2018) противопоставляют две модели, с
судьёй-человеком, решающим, какой аргумент более убедителен; надежда
состоит в том, что истину проще проверить, чем сгенерировать.
\emph{Рекурсивное моделирование вознаграждения} (Leike et al., 2018) использует
ИИ-помощника для помощи человеку в оценке, создавая иерархию
надзора. \emph{Constitutional AI} (Bai et al., 2022) заменяет
часть обратной связи от людей набором письменных принципов, которые модель
критикует сама против себя, снижая\,---\,хотя и не
устраняя,\,---\,нагрузку на оценщиков-людей.

\subsection{Фундаментальный вызов}

По своей сути задача выравнивания~--- это задача \emph{спецификации
ценностей}\index{value specification}: перевод богатого,
контекстно-зависимого и часто противоречивого ландшафта человеческих ценностей
в формальную цель. RL предоставляет оптимизационный механизм, но
не может сказать нам, \emph{что} оптимизировать. Это не просто
техническая задача; это философская, затрагивающая вопросы,
которые этики и политические теоретики обсуждают веками.

Практическое следствие для исследователей RL состоит в том, что проектирование вознаграждения
останется первоклассной заботой. Успехи в интерпретируемости,
лучшие методы получения предпочтений людей и формальные фреймворки
для спецификации ограничений~--- всё это сыграет свою роль. Никакая единственная техника
вряд ли «решит» выравнивание; вызов системный, и
прогресс будет постепенным.


% ===========================================================
\section{Безопасность и устойчивость}
\label{sec:open-safety}
% ===========================================================

\index{safe reinforcement learning}
\index{robustness}
Стратегия RL, хорошо работающая в среднем, может катастрофически отказывать
в редких, но важных ситуациях. Безопасность и устойчивость~---
не предметы роскоши; они являются предпосылками для развёртывания RL
в реальном мире.

\subsection{Безопасное исследование}
\index{safe exploration}
\index{constrained MDP}

В таких областях, как автономное вождение, здравоохранение и промышленное
управление, агент не может исследовать свободно: определённые состояния
(столкновения, передозировки, повреждение оборудования) должны избегаться даже во время
обучения. \emph{Ограниченные MDP}\index{constrained MDP} формализуют
это, дополняя стандартный MDP функциями стоимости и
ограничениями вида $\E_\pi[\sum_t c(s_t, a_t)] \leq d$.
Лагранжевы методы и ограниченная оптимизация стратегии (CPO;
Achiam et al., 2017) оптимизируют вознаграждение с учётом этих ограничений,
но гарантировать удовлетворение ограничений во время обучения\,---\,не
только при сходимости,\,---\,остаётся сложным.

\subsection{Распределительная устойчивость и состязательные атаки}
\index{distributional robustness}
\index{adversarial attack!on RL}

Стратегии, обученные в одной среде, могут сломаться при развёртывании в другой,
отличающейся даже незначительно. \emph{Распределительная устойчивость}
формулирует задачу как игру между агентом и противником,
возмущающим динамику переходов или наблюдения в пределах
ограниченного множества неопределённости. Робастные MDP (Iyengar, 2005; Nilim и
El~Ghaoui, 2005) обеспечивают теоретическую основу, и недавние работы
масштабировали эти идеи на глубокий RL через состязательное обучение и
рандомизацию домена.

Связанной озабоченностью являются \emph{состязательные атаки на стратегии RL}:
небольшие, тщательно подобранные возмущения наблюдения могут заставить
хорошо обученную стратегию предпринимать катастрофически плохие действия (Huang et
al., 2017). Сертифицированные защиты, аналогичные разработанным для
классификаторов изображений, для последовательного
принятия решений всё ещё находятся в зачаточном состоянии.

\subsection{Формальная верификация}
\index{formal verification}

Можем ли мы \emph{доказать}, что обученная стратегия удовлетворяет свойству безопасности?
Методы формальной верификации из программной инженерии\,---\,проверка моделей,
абстрактная интерпретация, выполнимость по модулю теорий
(SMT)\,---\,были применены к небольшим нейросетевым контроллерам, но
масштабирование их на сети, используемые в современном глубоком RL, остаётся открытой
задачей. Преодоление разрыва между статистическими гарантиями
теории обучения и детерминированными гарантиями формальных методов~---
одна из наиболее важных\,---\,и наиболее трудных,\,---\,задач в
области.


% ===========================================================
\section{Офлайн и пакетный RL}
\label{sec:open-offline}
% ===========================================================

\index{offline reinforcement learning}
\index{batch reinforcement learning}
Во многих практических обстановках\,---\,здравоохранение, рекомендательные системы,
диалог,\,---\,онлайн-взаимодействие дорогостояще или рискованно.
\emph{Офлайн RL} (также называемый \emph{пакетным RL}) стремится обучить
хорошую стратегию из фиксированного набора данных ранее собранных переходов
без дальнейшего взаимодействия со средой. Идея привлекательна:
организации часто располагают большими журналами прошлого поведения, и офлайн
RL обещает извлечь ценность из этих данных.

\subsection{Задача сдвига распределения}
\index{distribution shift!offline RL}

Центральная трудность~--- \emph{сдвиг распределения}. Набор данных был
собран некоторой поведенческой стратегией $\mu$, но обученная стратегия
$\pi$ может посещать состояния, которые $\mu$ редко или никогда не встречала. В
этих областях оценки ценности ненадёжны, и наивные офлайн-методы
могут расходиться или производить обманчиво оптимистичные стратегии.

\emph{Conservative Q-Learning} (CQL; Kumar et al., 2020) решает это,
добавляя член регуляризации, штрафующий Q-функцию на
внераспределённых действиях, эффективно обучая нижнюю границу истинной
ценности. Implicit Q-Learning (IQL; Kostrikov et al., 2022)
полностью избегает опроса стратегии во время обучения, используя ожидаемую
регрессию для оценки Q-значений верхней огибающей из набора данных.
Оба метода показали сильную производительность на стандартных тестах,
но их эффективность существенно зависит от качества и покрытия
набора данных.

\subsection{Моделирование последовательностей и Decision Transformer}
\index{Decision Transformer}

Влиятельной альтернативой офлайн RL на основе функции ценности является
переформулирование задачи как \emph{моделирования последовательностей}. Decision
Transformer (Chen et al., 2021) обусловливает трансформер на желаемом
доходе, прошлых состояниях и прошлых действиях, и авторегрессионно
предсказывает будущие действия. Во время тестирования агент обусловлен на
высоком доходе и генерирует соответствующее поведение. Этот подход
полностью обходит обновления Беллмана, используя сильные стороны трансформеров
в моделировании последовательностей.

Trajectory Transformer (Janner et al., 2021) расширил идею путём
моделирования совместного распределения по состояниям, действиям и вознаграждениям,
обеспечивая планирование через поиск по лучу во время вывода. Эти методы размывают
границу между RL и обучением с учителем, и их теоретические
свойства\,---\,когда они превосходят, когда отказывают и почему,\,---\,
всё ещё осмысляются.

\subsection{Переход от офлайн к онлайн}
\index{offline-to-online RL}

Многообещающим гибридом является \emph{переход от офлайн к онлайн}: предобучение стратегии
из статического набора данных с последующим дообучением небольшим объёмом онлайн-
взаимодействия. Это отражает парадигму предобучение-затем-дообучение, столь
успешную для языковых моделей
(Глава~\ref{ch:lm-post-training}). Ранние результаты (Nair et al.,
2020; Lee et al., 2022) показывают, что офлайн-предобучение может
резко ускорить онлайн-обучение, но предотвращение катастрофического
забывания офлайн-знаний и управление сдвигом распределения
при переходе остаются открытыми задачами.


% ===========================================================
\section{Мультимодальный и воплощённый RL}
\label{sec:open-embodied}
% ===========================================================

\index{embodied RL}
\index{sim-to-real transfer}
Обучение с подкреплением началось с физических интуиций\,---\,робот,
навигирующий по лабиринту, термостат, регулирующий температуру,\,---\,но за
последнее десятилетие наиболее заметные успехи были в симуляции:
играх, языке и рекомендациях. Возврат RL в физический
мир~--- один из великих незавершённых проектов области.

\subsection{Перенос из симуляции в реальность}

Обучение в симуляции дёшево и безопасно; развёртывание в реальном мире~---
ни то, ни другое. \emph{Разрыв симуляция-реальность}\index{sim-to-real gap}
возникает, потому что симуляторы неизбежно упрощают физику, шум датчиков
и визуальный вид. Рандомизация домена (Tobin et al., 2017)
смягчает разрыв, обучая стратегию по распределению симулированных сред, так
что реальный мир выглядит просто ещё одним образцом. Идентификация системы,
адаптивное управление и обученные остаточные модели предлагают дополнительные подходы,
но надёжное преодоление разрыва симуляция-реальность для
манипуляций с контактами остаётся активной задачей.

\subsection{Модели визуально-языково-действенные}
\index{vision-language-action model}

Новый класс \emph{визуально-языково-действенных} (VLA) моделей стремится
объединить восприятие, понимание языка и моторное управление в единой
архитектуре. RT-2 (Brohan et al., 2023) дообучил
визуально-языковую модель для вывода роботических действий как текстовых токенов,
демонстрируя, что семантические знания из крупномасштабного предобучения на
веб-данных могут переноситься на задачи манипуляции. Octo (Ghosh et al., 2024) и
$\pi_0$ (Black et al., 2024) продвинули это дальше с
архитектурами трансформеров, принимающими изображения и языковые инструкции
и выдающими непрерывные команды суставов.

Эти модели всё ещё далеки от ловкости и адаптируемости
человеческих манипуляций. Ключевые узкие места включают нехватку
высококачественных данных роботических демонстраций, сложность спецификации
вознаграждений для задач с контактами и вычислительные затраты
на вывод в реальном времени на аппаратуре с ограниченными ресурсами.

\subsection{Непрерывное управление в масштабе}

Помимо манипуляций, RL применяется к локомоции (шагающие
роботы, обученные через PPO в симуляции и развёрнутые на аппаратуре),
автономным гонкам (Gran Turismo), навигации дронов и управлению
гуманоидами всем телом. Общая нить состоит в том, что эти задачи включают
высокоразмерные непрерывные пространства действий, сложную динамику контактов
и жёсткие требования к реальному времени. Масштабирование RL на эти обстановки
раздвигает как алгоритмические, так и инженерные рубежи.


% ===========================================================
\section{RL и базовые модели}
\label{sec:open-foundation}
% ===========================================================

\index{foundation model!RL for}
Пожалуй, наиболее значимое развитие в недавнем RL~--- его появление
в качестве \emph{универсального метода дообучения} для базовых
моделей. Главы~\ref{ch:lm-post-training}--\ref{ch:reasoning}
зафиксировали это для текста; тот же паттерн теперь повторяется через
модальности.

\subsection{RL для языка и рассуждений}

RLHF (Глава~\ref{ch:dpo}) и GRPO (Глава~\ref{ch:grpo}) стали
стандартными инструментами для выравнивания и улучшения больших языковых моделей.
Как обсуждалось в Главе~\ref{ch:reasoning}, RL с верифицируемыми вознаграждениями
обеспечил кардинальные улучшения в математическом рассуждении и генерации кода.
Рубеж смещается в сторону задач рассуждений с более длинным горизонтом\,---\,многошаговое
научное решение задач, программная инженерия, доказательство теорем,\,---\,где
вознаграждение разрежено и отложено, а пространство действий~--- пространство
всех возможных продолжений текста.

\subsection{RL для визуальной генерации}

Тот же цикл RLHF, который улучшил языковые модели, был применён к
генерации изображений. DPOK (Fan et al., 2023) и DDPO (Black et al.,
2024) использовали методы градиента стратегии для дообучения диффузионных моделей
против моделей вознаграждения с предпочтениями людей, улучшая эстетическое качество и
соответствие промпту. Для генерации видео дообучение на основе RL использовалось
для улучшения временной согласованности и визуальной точности, хотя
вычислительные затраты на генерацию и оценку видеотраекторий остаются
узким местом.

\subsection{RL для научных открытий}
\index{scientific discovery!RL for}

RL внёс вклад в заметные научные прорывы. Обучающий конвейер AlphaFold
(Jumper et al., 2021) включал RL-подобное структурное уточнение. GNoME (Merchant et al., 2023) использовал RL-управляемый
поиск для открытия миллионов новых стабильных кристаллических структур.
FunSearch (Romera-Paredes et al., 2024) объединил большие языковые модели
с эволюционным поиском для открытия новых математических конструкций,
превзойдя людей-математиков в задаче о cap-set.

Общий паттерн~--- RL как \emph{внешний цикл}: базовая модель
предлагает кандидаты (структуры белков, составы материалов,
математические гипотезы), а RL-подобная цель оценивает и отбирает
среди них. Расширение этого паттерна на химию, разработку лекарств,
моделирование климата и инженерное проектирование~--- активный и
перспективный рубеж.


% ===========================================================
\section{Теоретические рубежи}
\label{sec:open-theory}
% ===========================================================

\index{RL theory}
Теория RL богата в табличных обстановках и обстановках с линейной аппроксимацией функций
(Главы~\ref{ch:mdp}--\ref{ch:td}), но методы, работающие
лучше всего на практике\,---\,глубокие сети, трансформеры, высокоразмерные
пространства стратегий,\,---\,часто опережают доступную теорию.

\subsection{Сожаление и сложность выборки для глубокого RL}
\index{regret bounds!deep RL}
\index{PAC-MDP}

Границы сожаления для табличных MDP тесные: алгоритмы, такие как UCBVI,
достигают $\tilde{O}(\sqrt{H^3 S A T})$ сожаления, где $H$~---
горизонт, $S$~--- число состояний, $A$~--- число действий, и $T$~---
число эпизодов. Для линейных MDP также известны совпадающие верхние и нижние
границы. Но для нелинейной аппроксимации функций,
используемой в глубоком RL\,---\,сверхпараметризованные нейронные сети, обученные с
стохастическим градиентным спуском,\,---\,сопоставимых гарантий не существует.

\emph{Анализ нейронного тангенциального ядра} (NTK) даёт некоторое понимание,
показывая, что достаточно широкие сети ведут себя приблизительно линейно,
но этот режим не захватывает обучение представлений, которое
делает глубокий RL эффективным. Разработка теории, учитывающей
обучение признаков, неконвексную оптимизацию и компромисс
исследование-использование одновременно,~--- одна из наиболее трудных открытых задач
в теории машинного обучения.

\subsection{Сходимость нелинейной аппроксимации функций}
\index{deadly triad}

\emph{Смертельная триада}\,---\,комбинация бутстрапа,
офлайн-обучения и аппроксимации функций,\,---\,была идентифицирована
Саттоном и Барто как источник расходимости в приближённом динамическом
программировании (Глава~\ref{ch:value-approx}). Хотя алгоритмы, такие
как DQN, хорошо работают на практике, доказательство их сходимости под
реалистичными допущениями (сети конечной ширины, неiid данные,
воспроизведение опыта) остаётся неуловимым. Недавний прогресс включает
результаты сходимости для сверхпараметризованных сетей в режиме NTK
и двухвременные анализы методов актор-критик, но единая
теория всё ещё отсутствует.

\subsection{Статистическая сложность RLHF}
\index{RLHF!statistical complexity}

Взрыв RLHF в обучении языковых моделей
(Глава~\ref{ch:dpo}) породил новые теоретические вопросы. Сколько
сравнений людей нужно для обучения модели вознаграждения, точно
отражающей человеческие предпочтения? При каких условиях цель DPO восстанавливает ту же стратегию, что конвейер RLHF? Как
качество данных предпочтений влияет на статистическую
эффективность выравнивания? Ранние результаты (Zhu et al., 2024; Xu et al.,
2024) устанавливают границы сложности выборки при
Bradley--Terry\index{Bradley--Terry model} модели предпочтений, но
теория пока не учитывает распределительные и стратегические
сложности реальной обратной связи от людей.

\subsection{Почему RL работает для LLM?}

Пожалуй, наиболее глубокая теоретическая загадка~--- почему дообучение RL работает
так хорошо для больших языковых моделей. Пространство действий огромно
(полный словарь на каждой позиции токена), вознаграждение
разрежено и отложено (часто единственный скаляр для целого ответа), и
стратегия~--- сеть с миллиардом параметров. Классическая теория RL
предсказывает, что такие задачи должны быть неразрешимы, однако PPO и GRPO
производят надёжные улучшения на практике. Понимание
этого\,---\,вероятно, через сочетание структуры естественного
языка, неявного исследования, обеспечиваемого сэмплированием токенов, и
тёплого старта от предобучения,\,---\,является ключевым открытым вопросом.


% ===========================================================
\section{Путь вперёд}
\label{sec:open-road}
% ===========================================================

Позвольте нам завершить, отступив от отдельных задач и рассмотрев
более широкую траекторию области.

\paragraph{RL как универсальное дообучение.}
Наиболее поразительная тенденция последних нескольких лет~--- обобщение RL
от нишевого метода для игр и робототехники
в \emph{универсальный метод дообучения} для любой модели, генерирующей
структурированный вывод. Языковые модели, генераторы изображений, видео-синтезаторы,
предсказатели структур белков и роботические контроллеры~---
все улучшаются через дообучение RL. Лежащий в основе
принцип прост: если вы можете оценить вывод, вы можете использовать RL для
того, чтобы модель производила выводы с более высокой оценкой. Этот принцип, впервые
формализованный в теореме о градиенте стратегии (Глава~\ref{ch:pg}),
оказывается одной из наиболее мощных идей в современном ИИ.

\paragraph{Сближение классической теории и современной практики.}
На протяжении десятилетий теория RL и практика RL развивались в
значительной мере отдельными сообществами. Теоретики изучали табличные MDP и линейную аппроксимацию
функций; практики обучали глубокие сети в масштабе.
Этот разрыв начинает сокращаться. Теория офлайн RL информирует практическое
проектирование алгоритмов; границы PAC-Bayes обеспечивают невакуумные гарантии обобщения
для обученных стратегий; а изучение RLHF приносит
инструменты теории обучения к одному из наиболее значимых
применений RL. Единая теория, соединяющая элегантные
результаты табличной обстановки с эмпирическими успехами глубокого RL,
стала бы знаковым достижением.

\paragraph{Агенты и автономия.}
Глава~\ref{ch:reasoning} и Глава~\ref{ch:agents} описали
появление агентов с обучением RL, способных рассуждать, использовать инструменты и
координироваться с другими агентами. Эта тенденция ускоряется. По мере того как
языковые модели становятся более мощными и дообучение RL обеспечивает
планирование с более длинным горизонтом, мы можем увидеть агентов, которые автономно проводят
научные исследования, пишут и отлаживают сложное программное обеспечение, управляют
инфраструктурой и ведут переговоры от имени своих пользователей. Технические
задачи огромны\,---\,отнесение заслуг через
тысячи шагов, безопасное исследование в открытых средах,
надёжная многоагентная координация,\,---\,но так же велики и потенциальные
выгоды.

\paragraph{Важность безопасности.}
С большей автономией приходит больший риск. Задачи безопасности и выравнивания,
обсуждённые в Разделах~\ref{sec:open-alignment}
и~\ref{sec:open-safety}, будут только усугубляться по мере того, как агенты
становятся более мощными. Области нужны не только лучшие алгоритмы, но и
лучшие институты: стандарты оценки, практики красных команд,
регуляторные рамки и культура ответственного
развития. Исследователи RL несут здесь особую ответственность, поскольку
оптимизационное давление, оказываемое RL,~--- одно из наиболее сильных во
всём машинном обучении.

\paragraph{Что может принести следующее десятилетие.}
Предсказывать будущее быстро движущейся области~--- занятие дурака,
но несколько тенденций, по всей видимости, сохранятся. Во-первых, масштаб
обучения RL продолжит расти, движимый экономикой
дообучения базовых моделей. Во-вторых, граница между RL,
обучением с учителем и самосупервизированным обучением продолжит
размываться, как демонстрируют методы, такие как DPO (Глава~\ref{ch:dpo}) и Decision
Transformer (Раздел~\ref{sec:open-offline}). В-третьих, RL
будет применяться ко всё более широкому спектру областей: от науки о
материалах до климатической политики и образования. В-четвёртых, задачи
безопасности и выравнивания вынудят сообщество разработать новые инструменты для
оценки, надзора и управления.

История обучения с подкреплением далека от завершения. Если что, она
только достигает своей наиболее интересной главы. Математические
основы, заложенные в Частях~I и~II этой книги, алгоритмический
инструментарий, разработанный в Части~III, и передовые приложения, обозреваемые
в Части~IV,~--- всё это обеспечивает словарь, необходимый для внесения вклада. Открытые
задачи, описанные в этой главе,~--- не препятствия для сожаления;
они~--- возможности для использования. Мы надеемся, что эта книга вооружила вас
для их использования.
